% GENERATED FILE. Edit canonical lessons, then run npm run generate:books.
% canonical-source-hash: fnv1a64:c7eb50345e31b779
% canonical-lessons: TE-C384-sucincu, TE-C384-pratipadincu, TE-C384-pratipadana, TE-C384-sipharsu-ceyu, TE-C384-amsam

\chapter{Suggestions and Proposals}
\label{ch:te-suggestions-and-proposals}
\hlchaptermodality{384}

\begin{chapteropening}
\textbf{By the end of this chapter:} \emph{I can suggest, propose and recommend.}

Complete the last lesson of chapter 384: I can suggest, propose and recommend.
\end{chapteropening}

\section[\texorpdfstring{\te{సూచించు} (sūciñcu)}{సూచించు (sūciñcu)}]{\te{సూచించు} (sūciñcu) — to suggest; to indicate}

\label{lesson:TE-C384-sucincu}



\begin{quote}
Before the new one: say the Telugu for logic, then the Telugu for reasonable.
\end{quote}

\subsection*{You'll want to know: \te{సూచించు}}
\textbf{\te{సూచించు}} — \emph{sūciñcu} — \textquotedblleft{}to suggest; to indicate\textquotedblright{}.

\begin{quote}
\te{మా} \te{ఉపాధ్యాయుడు} \te{ఒక} \te{కొత్త} \te{పుస్తకం} \te{సూచించాడు}.
\end{quote}

\begin{quote}
\emph{mā upādhyāyuḍu oka kotta pustakaṁ sūciñcāḍu.}
\end{quote}

\begin{quote}
Our teacher suggested a new book.
\end{quote}

Past (I …): \te{సూచించాను}, \emph{sūciñcānu}.

The verb of \te{సూచన}, a suggestion, which you know.

\begin{cousinweb}
\textbf{\te{సూచించు}} (\emph{sūciñcu}) is built on Sanskrit \emph{sūcana}, \textquotedblleft{}pointing out, indicating; an indication\textquotedblright{} (Monier-Williams), with the Telugu verb ending -\te{ఇంచు}.
\end{cousinweb}

\begin{grammarlens}[title={What you've built}]
One more word for giving reasons and opinions.
\end{grammarlens}

\subsection*{Your turn}
\begin{itemize}
\raggedright
  \item \emph{Say it:} \emph{sūciñcu}
  \item \emph{Say it:} \emph{sūciñcu}, once more
  \item \emph{Say it:} \emph{mā upādhyāyuḍu oka kotta pustakaṁ sūciñcāḍu.} — Our teacher suggested a new book.
  \item \emph{Recall:} say the Telugu for logic, then the Telugu for reasonable, then say \emph{sūciñcu} again
\end{itemize}

\begin{tcolorbox}[breakable,colback=teal!4,colframe=teal!35!black,title={Before you move on}]
What does \te{సూచించు} mean? (\textquotedblleft{}To suggest; to indicate\textquotedblright{}.) Say it once more.
\end{tcolorbox}
\section[\texorpdfstring{\te{ప్రతిపాదించు} (pratipādiñcu)}{ప్రతిపాదించు (pratipādiñcu)}]{\te{ప్రతిపాదించు} (pratipādiñcu) — to propose}

\label{lesson:TE-C384-pratipadincu}



\begin{quote}
Before the new one: say the Telugu for reasonable, then the Telugu for to suggest.
\end{quote}

\subsection*{You'll want to know: \te{ప్రతిపాదించు}}
\textbf{\te{ప్రతిపాదించు}} — \emph{pratipādiñcu} — \textquotedblleft{}to propose\textquotedblright{}.

\begin{quote}
\te{ఆమె} \te{సమావేశంలో} \te{ఒక} \te{కొత్త} \te{పద్ధతి} \te{ప్రతిపాదించింది}.
\end{quote}

\begin{quote}
\emph{āme samāvēśaṁlō oka kotta paddhati pratipādiñcindi.}
\end{quote}

\begin{quote}
She proposed a new method at the meeting.
\end{quote}

Past (I …): \te{ప్రతిపాదించాను}, \emph{pratipādiñcānu}.

Formal: meetings, committees and papers. In everyday talk \te{సూచించు} does the work.

\begin{cousinweb}
\textbf{\te{ప్రతిపాదించు}} (\emph{pratipādiñcu}) is built on Sanskrit \emph{pratipādana}, \textquotedblleft{}giving, presenting\textquotedblright{} (Monier-Williams), with the Telugu verb ending -\te{ఇంచు}; Telugu uses it for putting an idea forward.
\end{cousinweb}

\begin{grammarlens}[title={What you've built}]
One more word for giving reasons and opinions.
\end{grammarlens}

\subsection*{Your turn}
\begin{itemize}
\raggedright
  \item \emph{Say it:} \emph{pratipādiñcu}
  \item \emph{Say it:} \emph{pratipādiñcu}, once more
  \item \emph{Say it:} \emph{āme samāvēśaṁlō oka kotta paddhati pratipādiñcindi.} — She proposed a new method at the meeting.
  \item \emph{Recall:} say the Telugu for reasonable, then the Telugu for to suggest, then say \emph{pratipādiñcu} again
\end{itemize}

\begin{tcolorbox}[breakable,colback=teal!4,colframe=teal!35!black,title={Before you move on}]
What does \te{ప్రతిపాదించు} mean? (\textquotedblleft{}To propose\textquotedblright{}.) Say it once more.
\end{tcolorbox}
\section[\texorpdfstring{\te{ప్రతిపాదన} (pratipādana)}{ప్రతిపాదన (pratipādana)}]{\te{ప్రతిపాదన} (pratipādana) — a proposal}

\label{lesson:TE-C384-pratipadana}



\begin{quote}
Before the new one: say the Telugu for to suggest, then the Telugu for to propose.
\end{quote}

\subsection*{You'll want to know: \te{ప్రతిపాదన}}
\textbf{\te{ప్రతిపాదన}} — \emph{pratipādana} — \textquotedblleft{}a proposal\textquotedblright{}.

\begin{quote}
\te{మా} \te{ప్రతిపాదనను} \te{అందరూ} \te{ఒప్పుకున్నారు}.
\end{quote}

\begin{quote}
\emph{mā pratipādananu andarū oppukunnāru.}
\end{quote}

\begin{quote}
Everyone accepted our proposal.
\end{quote}

Plural: \te{ప్రతిపాదనలు}, \emph{pratipādanalu}.

The noun of \te{ప్రతిపాదించు}: \te{ఒక} \te{ప్రతిపాదన} \te{చేయు}, to make a proposal.

\begin{cousinweb}
\textbf{\te{ప్రతిపాదన}} (\emph{pratipādana}) is borrowed from Sanskrit \emph{pratipādana}, \textquotedblleft{}giving, presenting\textquotedblright{} (Monier-Williams); its verb is \te{ప్రతిపాదించు}.
\end{cousinweb}

\begin{grammarlens}[title={What you've built}]
One more word for giving reasons and opinions.
\end{grammarlens}

\subsection*{Your turn}
\begin{itemize}
\raggedright
  \item \emph{Say it:} \emph{pratipādana}
  \item \emph{Say it:} \emph{pratipādana}, once more
  \item \emph{Say it:} \emph{mā pratipādananu andarū oppukunnāru.} — Everyone accepted our proposal.
  \item \emph{Recall:} say the Telugu for to suggest, then the Telugu for to propose, then say \emph{pratipādana} again
\end{itemize}

\begin{tcolorbox}[breakable,colback=teal!4,colframe=teal!35!black,title={Before you move on}]
What does \te{ప్రతిపాదన} mean? (\textquotedblleft{}A proposal\textquotedblright{}.) Say it once more.
\end{tcolorbox}
\section[\texorpdfstring{\te{సిఫార్సు} \te{చేయు} (siphārsu cēyu)}{సిఫార్సు చేయు (siphārsu cēyu)}]{\te{సిఫార్సు} \te{చేయు} (siphārsu cēyu) — to recommend}

\label{lesson:TE-C384-sipharsu-ceyu}



\begin{quote}
Before the new one: say the Telugu for to propose, then the Telugu for a proposal.
\end{quote}

\subsection*{You'll want to know: \te{సిఫార్సు} \te{చేయు}}
\textbf{\te{సిఫార్సు} \te{చేయు}} — \emph{siphārsu cēyu} — \textquotedblleft{}to recommend\textquotedblright{}.

\begin{quote}
\te{నా} \te{స్నేహితుడు} \te{ఆ} \te{వైద్యుడిని} \te{సిఫార్సు} \te{చేశాడు}.
\end{quote}

\begin{quote}
\emph{nā snēhituḍu ā vaidyuḍini siphārsu cēśāḍu.}
\end{quote}

\begin{quote}
My friend recommended that doctor.
\end{quote}

Past (I …): \te{సిఫార్సు} \te{చేశాను}, \emph{siphārsu cēśānu}.

\te{సిఫార్సు} alone is a recommendation, and also a word put in for someone.

\begin{cousinweb}
\textbf{\te{సిఫార్సు} \te{చేయు}} (\emph{siphārsu cēyu}) is a loan from Persian \emph{sifārish}, \textquotedblleft{}a recommendation\textquotedblright{}, with the native verb \te{చేయు}, \textquotedblleft{}to do\textquotedblright{} (DEDR 1957).
\end{cousinweb}

\begin{grammarlens}[title={What you've built}]
One more word for giving reasons and opinions.
\end{grammarlens}

\subsection*{Your turn}
\begin{itemize}
\raggedright
  \item \emph{Say it:} \emph{siphārsu cēyu}
  \item \emph{Say it:} \emph{siphārsu cēyu}, once more
  \item \emph{Say it:} \emph{nā snēhituḍu ā vaidyuḍini siphārsu cēśāḍu.} — My friend recommended that doctor.
  \item \emph{Recall:} say the Telugu for to propose, then the Telugu for a proposal, then say \emph{siphārsu cēyu} again
\end{itemize}

\begin{tcolorbox}[breakable,colback=teal!4,colframe=teal!35!black,title={Before you move on}]
What does \te{సిఫార్సు} \te{చేయు} mean? (\textquotedblleft{}To recommend\textquotedblright{}.) Say it once more.
\end{tcolorbox}
\section[\texorpdfstring{\te{అంశం} (aṁśaṁ)}{అంశం (aṁśaṁ)}]{\te{అంశం} (aṁśaṁ) — a point; an issue}

\label{lesson:TE-C384-amsam}



\begin{quote}
Before the new one: say the Telugu for a proposal, then the Telugu for to recommend.
\end{quote}

\subsection*{You'll want to know: \te{అంశం}}
\textbf{\te{అంశం}} — \emph{aṁśaṁ} — \textquotedblleft{}a point; an issue\textquotedblright{}.

\begin{quote}
\te{సమావేశంలో} \te{మూడు} \te{ముఖ్యమైన} \te{అంశాలు} \te{చర్చించారు}.
\end{quote}

\begin{quote}
\emph{samāvēśaṁlō mūḍu mukhyamaina aṁśālu carciñcāru.}
\end{quote}

\begin{quote}
They discussed three important points at the meeting.
\end{quote}

Plural: \te{అంశాలు}, \emph{aṁśālu}.

In the plural the -\te{ం} drops and -\te{ాలు} comes in, as \te{పుస్తకం} becomes \te{పుస్తకాలు}.

\begin{cousinweb}
\textbf{\te{అంశం}} (\emph{aṁśaṁ}) is borrowed from Sanskrit \emph{aṁśa}, \textquotedblleft{}a share, a part\textquotedblright{} (CDIAL 2). In Telugu a part of a discussion is a point or an issue.
\end{cousinweb}

\begin{grammarlens}[title={What you've built}]
One more word for giving reasons and opinions.
\end{grammarlens}

\subsection*{Your turn}
\begin{itemize}
\raggedright
  \item \emph{Say it:} \emph{aṁśaṁ}
  \item \emph{Say it:} \emph{aṁśaṁ}, once more
  \item \emph{Say it:} \emph{samāvēśaṁlō mūḍu mukhyamaina aṁśālu carciñcāru.} — They discussed three important points at the meeting.
  \item \emph{Recall:} say the Telugu for a proposal, then the Telugu for to recommend, then say \emph{aṁśaṁ} again
\end{itemize}

\begin{tcolorbox}[breakable,colback=teal!4,colframe=teal!35!black,title={Before you move on}]
What does \te{అంశం} mean? (\textquotedblleft{}A point; an issue\textquotedblright{}.) Say it once more.
\end{tcolorbox}
